% Work-MAT | 2. rocnik | pripravovana hodina 60
% Návaznost: uživatelské připravy.zip (hodiny 01–20), autorská učebnice na GitHubu.
% Realisticky.cz: didaktická inspirace v posloupnosti a typech úloh; zadání jsou vlastní.
\documentclass[11pt,a4paper]{article}
\input{../../sablony/pripravy/priprava.sty}
\usepackage{tikz}
\nastavPripravu{2. ročník}{Goniometrické vzorce}{60}{Vzorce pro dvojnásobný úhel}
\begin{document}
\blok{Motivace a návaznost}
Dvojnásobný úhel odvodíme ze součtových vzorců dosazením $\beta=\alpha$. Při počítání si žáci připomenou kvadranty.
\poznamka{Orientační nabídka pro 45--55 minut: krátké zopakování a motivace, nový výklad, tři společné vzorové úlohy, samostatná práce. Vyber podle tempa třídy; zbývající zadání slouží jako rezerva.}
\blok{Výklad}
$\sin2x=2\sin x\cos x$ a $\cos2x=\cos^2x-\sin^2x=1-2\sin^2x=2\cos^2x-1$.
\blok{Společné příklady}
\begin{prikladReseniBox}[title={Příklad 1}]
Vyřeš, resp. urči: $\sin60^\circ\text{ ze vzorce pro dvojnásobek}$.
\tcblower
\[\begin{aligned}2\sin30^\circ\cos30^\circ=2\cdot\frac{1}{2}\cdot\frac{\sqrt3}{2} \\ =\frac{\sqrt3}{2}\end{aligned}\]
\end{prikladReseniBox}
\begin{prikladReseniBox}[title={Příklad 2}]
Vyřeš, resp. urči: $\cos120^\circ\text{ ze vzorce pro dvojnásobek}$.
\tcblower
\[\begin{aligned}1-2\sin^260^\circ=1-2\cdot\frac{3}{4} \\ =-\frac{1}{2}\end{aligned}\]
\end{prikladReseniBox}
\poznamka{Pro rychlejší žáky přidej druhé odvození kosinového vzorce; hlavní třída musí zvládnout použití jedné správné varianty.}
\blok{Samostatné procvičování}
\begin{zadaniReseniBox}[title={Samostatná práce 1}]
\begin{enumerate}[label=\alph*)]
\item $1-\sin^2x$.
\item $1-\cos^2x$.
\item $\sin^2x+\cos^2x$.
\item $\frac{1-\cos^2x}{\sin x}$.
\item $\frac{1-\sin^2x}{\cos x}$.
\item $\operatorname{tg}x\cdot\cos x$.
\end{enumerate}
\tcblower
\begin{enumerate}[label=\alph*)]
\item $\cos^2x$.
\item $\sin^2x$.
\item $1$.
\item $\sin x\;(\sin x\ne0)$.
\item $\cos x\;(\cos x\ne0)$.
\item $\sin x\;(\cos x\ne0)$.
\end{enumerate}
\end{zadaniReseniBox}
\blok{Dobrovolné rozšíření}
\begin{zadaniReseniBox}[title={Rozšíření 1}]
\begin{enumerate}[label=\alph*)]
\item Navrhni obdobnou úlohu s jiným úhlem a zkontroluj výsledek graficky..
\end{enumerate}
\tcblower
\begin{enumerate}[label=\alph*)]
\item $Srovnej s jednotkovou kružnicí nebo grafem.$.
\end{enumerate}
\end{zadaniReseniBox}
\blok{Úkol na doma}
U každé úpravy napiš použitý vzorec a vysvětli případné podmínky, za kterých je vyjádření definované.
\begin{zadaniReseniBox}[title={Domácí úkol}]
\begin{enumerate}[label=\alph*)]
\item $\operatorname{cotg}x\cdot\sin x$.
\item $\sin(2x)$.
\item $\cos(2x)$.
\item $\cos(2x)$.
\item $\sin(30^\circ+60^\circ)$.
\item $\cos(45^\circ+45^\circ)$.
\item $\sin(45^\circ-45^\circ)$.
\item $\cos(60^\circ-30^\circ)$.
\end{enumerate}
\tcblower
\begin{enumerate}[label=\alph*)]
\item $\cos x\;(\sin x\ne0)$.
\item $2\sin x\cos x$.
\item $1-2\sin^2x$.
\item $2\cos^2x-1$.
\item $1$.
\item $0$.
\item $0$.
\item $\frac{\sqrt3}{2}$.
\end{enumerate}
\end{zadaniReseniBox}
\poznamka{Domácí úkol obsahuje 8 položek; předpokládaná domácí práce přibližně 15--30 minut. Po větších blocích jsou vědomě promíchané dříve probrané typy.}
\end{document}
